\documentclass{article}
\usepackage[T1]{fontenc}
\setlength{\parindent}{0pt}

\usepackage[margin=2.5cm]{geometry}

\usepackage{xcolor}
\usepackage{listings}
\lstset{
	basicstyle=\ttfamily,
	showstringspaces=false,
	keywordstyle=\color{blue},
	commentstyle=\color{gray},
	stringstyle=\color{violet},
	identifierstyle=,
	% numbers=left,
	% numberstyle=\scriptsize
}

\newcommand{\VE}[1]{\begin{verbatim}#1\end{verbatim}}

\title{\textbf{Some notes on slurm}}
\author{From ``Introduction to Slurm'', by SchedMD Slurm}
\date{\today}



\begin{document}

\maketitle

\section{Daemons}
I think we don't really need thease but I report them
\begin{itemize}
	\item \verb|slurmctld| - central controller (typically one per
		cluster).\\[3pt]This guy monitors everything, resources, job ques,
		and \textbf{allocates resources}.
	\item \verb|slurmd| - compute node daemon (typically one per compute
		node).\\[3pt]Launches and manages tasks.
	\item \verb|slurmdbd| - database daemon (typically one per
		enterprise).\\[3pt] Something something backup.
\end{itemize}

\section{Resource allocation commands}
I think we need this - remember that we always have \verb|man|!
Some things to keep in mind
\begin{itemize}
	\item almost all options have to formats: single letter (\verb|-p debug|)
		and verbose (\verb|--partition=debug|)
	\item time formats are \verb|days-hours:minutes:seconds|
	\item almost all commands support logging with \verb|-v| (more
		\verb|v|s-s for more verbosity)
	\item environment variables can be used to establish site/user-specific
		defaults (e.g. \verb|SQUEUE_STATES=all| for \verb|squeue| to display
		jobs in any state)
\end{itemize}

\medskip Ok now to the commands:
\begin{itemize}
	\item \verb|sbatch| - submit script for later execution

\begin{lstlisting}[language=bash, frame=single]
# Job 1
$ sbatch --ntasks=1 --time=10 preprocess.sh
Submitted batch job 45001

# Job 2 dependent on job 1
$ sbatch --ntasks=128 --time=60 --depend=45001 do_work.sh
Submitted batch job 45002

# Job 3 dependent on job 2
$ sbatch --ntasks=1 --time=30 --depend=45002 postprocess.sh
Submitted batch job 45003
\end{lstlisting}

	\item \verb|salloc| - create job allocatio nand start a shell to use it
		(interactive mode)

\begin{lstlisting}[language=bash, frame=single]
# Create allocation for 8 tasks and 10 minutes for bash shell, then
# launch some tasks
$ salloc --ntasks=8 --time=10 bash
	> salloc: Granted job allocation 45000
# env | grep SLURM
	> ...
\end{lstlisting}

	\item \verb|srun| - create a job allocatoin (if needed) and launch a job
		step (typically an MPI job)

\begin{lstlisting}[language=bash, frame=single]
# Create allocation for 2 tasks then launch "hostname" on the 
# allocation, label the output with the task ID
$ srun --ntasks=2 --label=hostname
0: tux123
1: tux123 

# As above, but allocate the job two whole nodes
$ srun --nnodes=2 --exclusive --label hostname
0: tux123
1: tux123 
\end{lstlisting}

		Interestingly, you can run different executables specifying different
		Task IDs using a configuration file and specifying
		\verb|--multi-prog <CONF_FILE>| 

\begin{lstlisting}[language=bash, frame=single]
$ cat master.conf
# TaskID	Program		Arguments
0		/usr/me/master	
1-4		/usr/me/slave	--rank=%0
$ srun --ntasks=5 --multi-prog master.conf
\end{lstlisting}

	\item \verb|sattach| - connect stdin/out/err for an existing job or job step
\end{itemize}

\section{State information commands}
\begin{itemize}
	\item \verb|sinfo| - reports system status (nodes, queues, ...).
		You can format the output how you want, blabla.

\begin{lstlisting}[language=bash, frame=single]
# Report status in node-oriented form
$ sinfo --Node 

# Report status of notdes in partition 'debug'
$ sinfo -p debug

# Report status every 60 seconds
$ squeue -i60
\end{lstlisting}

	\item \verb|squeue| - reports job and job step status - I think this is the
		one we will be using most of the time

\begin{lstlisting}[language=bash, frame=single]
# Report jobs for the user 'alec' in any state
$ squeue -u alec -t all

# Report steps in partition 'debug')
$ squeue -s -p debug

# Report currently active jobs every 60 seconds)
$ squeue -i60
\end{lstlisting}

	\item \verb|smap| - report system, job, or step status with topology
	\item \verb|sview| - newer version of \verb|smap|, report and/or update
		system, job, step partition or reservation status with topology
		(GTK-based GUI!)
	\item \verb|scontrol| - administrator tool, we are not interested
\end{itemize}


\section{Submitting jobs}
There are default parameters to submit a job:

\begin{itemize}
\item \verb|-c <NUM CPUS>|\\Number of CPUs that you want for your task
\item \verb|--mem-per-cpu <mem>|\\How much RAM we want per CPU (e.g. 100M)
\item \verb|-t <time>|\\How long you need the resources - the higher the time
	the lower the priority! (e.g. 3-0:0:0)
\item \verb|-J <JOB NAME>|\\Name of the job to run
\end{itemize}

There are also some specific ones:
\begin{itemize}
	\item \verb|--wrap <COMMAND>|\\Execute the command in a non blocking manner!
\end{itemize}

\begin{lstlisting}[language=bash, caption={sbatch: jobs in the
background}, frame=single, firstnumber=1]
$ sbatch -c 2 --mem-per-cpu 2G -t 5:0:0 \\
	 -J crunchy --wrap "python crunchy.py"
\end{lstlisting}

\begin{lstlisting}[language=bash, caption={srun: jobs in the
foreground}, frame=single]
$ sbatch -c 2 --mem-per-cpu 2G -t 5:0:0 \\
	 -J crunchy --wrap "python crunchy.py"
\end{lstlisting}

\begin{lstlisting}[language=bash, caption={salloc: allocate and do
nothing}, frame=single]
$ sbatch -c 2 --mem-per-cpu 2G -t 5:0:0 \\
	 -J crunchy 
\end{lstlisting}
\end{document}
